\begin{table}[H]
\centering
\caption{Standardized Effect Sizes for Main Outcomes}
\label{tab:sde}
\begin{threeparttable}
\small
\begin{adjustbox}{max width=\textwidth}
\begin{tabular}{llccccc}
\toprule
Outcome & Specification & $\hat{\beta}$ & SD($Y$) & SDE & SE(SDE) & Classification \\
\midrule
\multicolumn{7}{l}{\textit{Panel A: Pooled}} \\
Below 10-ton & Baseline & 0.0042 & 0.434 & 0.0097 & 0.0298 & Small positive \\
\midrule
\multicolumn{7}{l}{\textit{Panel B: Heterogeneous}} \\
Below 10-ton & Manufacturing & -0.0170 & 0.452 & -0.0376 & 0.0408 & Small negative \\
Below 10-ton & Non-Manufacturing & 0.0273 & 0.410 & 0.0666 & 0.0474 & Moderate positive \\
\bottomrule
\end{tabular}
\end{adjustbox}
\begin{tablenotes}[flushleft]
\small
\item \textit{Notes:} \textbf{Country:} United States. \textbf{Research question:} Does the removal of the ``Once In Always In'' regulatory constraint cause industrial facilities to strategically reduce reported hazardous air pollutant emissions below the Clean Air Act major source threshold? \textbf{Policy mechanism:} The 2018 EPA withdrawal of the OIAI guidance removed the rule that facilities classified as major HAP sources must permanently remain under costly MACT standards, creating a new escape hatch that allows major sources to reclassify as area sources by reducing any single HAP below 10 tons per year. \textbf{Outcome definition:} Binary indicator equal to one if a facility's highest individual HAP emission is below the 10-ton-per-year major source threshold in the NEI annual facility summary. \textbf{Treatment:} Binary; equal to one for all facility-years after January 25, 2018 (OIAI withdrawal date). \textbf{Data:} EPA National Emissions Inventory Facility Summaries, 2012--2021, facility-by-year panel. \textbf{Method:} Two-way fixed effects (facility and year FE), state-clustered standard errors. \textbf{Sample:} Facilities with max single-HAP emissions between 3 and 25 tons per year (near-threshold sample). SDE $= \hat{\beta} / \text{SD}(Y)$ where SD($Y$) is the pre-treatment standard deviation. Classification refers to magnitude, not statistical significance: Large ($|$SDE$| > 0.15$), Moderate ($0.05$--$0.15$), Small ($0.005$--$0.05$), Null ($< 0.005$). Classification labels refer to the magnitude of the standardized point estimate, not to statistical significance. ``Null'' denotes a near-zero effect size ($|$SDE$| < 0.005$), not a failure to reject a null hypothesis.
\end{tablenotes}
\end{threeparttable}
\end{table}

